\appendix

\section{The Causal Inference Toolkit: Historical Development and Parallels with Other Disciplines}
\label{sec:appendix-history}

This appendix provides a detailed account of the intellectual traditions that produced the causal inference toolkit introduced in Section~\ref{sec:bg-causal-history}, followed by a summary of methodological reforms in psychology and epidemiology in which they recently started to address causal infernece problems similar to those we identify in SE\@.

\subsection{The Development of the Causal Inference Toolkit}
\label{sec:appendix-toolkit-history}

The modern causal inference toolkit emerged over a century through contributions from three intellectual traditions---statistics, econometrics, and computer science---each addressing a distinct aspect of causal inference: experimental foundations, credibility evolution, and graphical modeling.

First, statisticians laid the \emph{experimental foundations}.
\citet{neyman1923application} defined the causal effect as the difference between two potential outcomes---only one of which is ever observed---and \citet{fisher1935design} established randomization as the principled basis for causal attribution.
Extending causal reasoning to observational settings, \citet{hill1965environment} proposed viewpoints for evaluating whether observed associations reflect genuine causal relationships, \citet{rubin1974estimating} formalized \emph{conditional ignorability} as the identifying assumption for observational studies, and \citet{rosenbaum1983central} introduced the propensity score as a practical dimension-reducing device for covariate adjustment.
Yet, the family of selection-on-observables methods fundamentally cannot address confounding from \emph{unobserved variables} plaguing many economic and societal problems~\cite{heckman1979sample}---e.g., one's (arguably unmeasurable) inherent ability affects both their educational attainment and future earnings~\cite{card1999causal}---a limitation that motivated economists to develop design-based alternatives.

Then, econometricians drove the \emph{credibility revolution}, shifting the basis of causal claims from statistical modeling to research design.
\citet{leamer1983let} first diagnosed the field's credibility problem, arguing that the sensitivity of regression results to specification choices undermined confidence in empirical findings.
The constructive response came through landmark natural experiments: \citeauthor{card1994minimum}'s~\cite{card1994minimum} difference-in-differences analysis of a minimum wage increase demonstrated that design-based evidence can overturn decades of cross-sectional conventional wisdom, and \citet{angrist1996identification} formalized instrumental variable identification through the LATE theorem.
\citet{angrist2010credibility} crystallized the overarching manifesto: Design-based studies yield more credible causal evidence than advances in statistical modeling alone.
The 2021 Nobel Prize in Economic Sciences, awarded to Card, Angrist, and Imbens for \emph{``their methodological contributions to the analysis of causal relationships''}~\cite{nobelprize2021economics}, cemented the credibility revolution as one of the most consequential intellectual movements in modern social science research.

From computer science and artificial intelligence (specifically the line of research on Bayesian networks and probabilistic reasoning~\cite{pearl1988probabilistic}), \citet{pearl2009causality} developed the graphical causal model framework, representing causal relationships as directed acyclic graphs (DAGs) and introducing the \emph{do-calculus} and \emph{back-door criterion} for determining sufficient adjustment sets.
Where the potential outcomes tradition defines \emph{what} to estimate, DAGs provide a visual language for arguing \emph{why} identification assumptions hold.
\citet{morgan2015counterfactuals} bridged the two traditions for social science, and \citet{imbens2020potential} clarified their relative strengths, arguing that potential outcomes are more natural for design-based research while DAGs excel at encoding complex causal structures.

Recent cross-disciplinary advances have extended the toolkit further: double/debiased machine learning for high-dimensional settings~\cite{chernozhukov2018double}, causal forests for heterogeneous effects~\cite{athey2016recursive}, modern DiD estimators robust to staggered adoption~\cite{callaway2021difference, roth2023whats}, and sensitivity analysis tools for assessing vulnerability to unobserved confounding~\cite{rosenbaum2002observational, oster2019unobservable, cinelli2020making}.
A new generation of accessible textbooks~\cite{angrist2009mostly, hernan2020causal, cunningham2021causal, huntingtonklein2021effect} has demonstrated that these methods can be taught to non-specialist audiences; \citeauthor{hernan2020causal}'s~\cite{hernan2020causal} \emph{target trial} framework, which asks researchers to describe the hypothetical randomized trial their observational analysis attempts to emulate, is particularly relevant to SE\@.
Our tutorial treats both traditions as essential for SE research: Design-based identification provides credible effect estimates for SE interventions while graphical causal models help the decomposition of mechanisms in which the intervention affects the outcome---both are critical for a full understanding of empirical SE problems (Section~\ref{sec:primer}).


\subsection{The Methodological Reform in Psychology and Epidemiology}
\label{sec:appendix-psych-reform}

In recent years, psychology and health research have confronted a causal inference problem similar to what we identify in SE---the disconnect between causal ambition and methodological practice---and their experience reveals both consequences and remedies relevant to the SE community.

The reform began with the diagnosis of the causal language problem.
\citet{hernan2018cword} challenged the epidemiological norm of replacing causal language with associational euphemisms, arguing that this evasion does not prevent causal interpretation but merely obscures the causal question and blocks transparent discussion of identification assumptions.
\citet{grosz2020taboo} documented a parallel ``taboo'' in psychology: The norm against explicit causal inference pushes causal reasoning underground, resulting in impaired study design and a disconnect between findings and downstream interpretation.
\citet{haber2022causal} provided large-scale empirical evidence that this taboo fails on its own terms---among 1,170 articles from 18 high-profile health journals, 52\% of abstracts implied moderate or strong causality despite ostensibly associational language.
The SE community exhibits the same pattern, as we will show in Section~\ref{sec:bg-adoption}.

In response, methodological reformers borrowed tools from the causal inference traditions described above.
\citet{rohrer2018thinking} introduced DAGs to psychology, and \citet{wysocki2022statistical} demonstrated that controlling for mediators, colliders, or their descendants can introduce \emph{more} bias---a finding directly applicable to SE studies that include covariates without causal justification.
\citet{lundberg2021estimand} formalized the \emph{estimand-first} principle for sociology, requiring every study to define its target quantity outside any particular statistical model.
Institutionally, \citet{lederer2019control} issued editor-endorsed guidance requiring authors to specify causal questions and justify covariate selection---editorial policy that does not yet exist in any SE venue.
\citet{tennant2021dags} assessed DAG adoption in health research and found that only 21\% of articles reported target estimands, revealing a theory-practice gap that parallels \citeauthor{DBLP:journals/infsof/Siebert23}'s~\cite{DBLP:journals/infsof/Siebert23} findings in SE\@.

High-profile case studies have illustrated the consequences of ignoring causal reasoning.
\citet{killingsworth2023income} appeared to resolve the long-standing income--happiness debate through sophisticated observational analysis, only for \citet{rohrer2024inappropriate} to demonstrate that the resolution rested on untested identification assumptions---effectively performing the causal assessment our framework codifies.
This sequence mirrors the trajectory from \citet{DBLP:conf/sigsoft/RayPFD14} through \citeauthor{DBLP:journals/toplas/BergerHMVV19}'s~\cite{DBLP:journals/toplas/BergerHMVV19} in SE: In both cases, increasingly sophisticated correlational analyses cannot resolve causal questions without explicit identification.
\citeauthor{rohrer2024causal}'s~\cite{rohrer2024causal} recent tutorial for psychologists is the closest existing analog to our work but does not cover design-based identification strategies or SE-specific confounding structures.
The solutions emerging across these disciplines---explicit causal questions, DAG-based covariate justification, estimand-first workflows, sensitivity analysis, and editorial enforcement---can inform and accelerate SE's own methodological reform.


\section{Frequently Asked Questions}
\label{sec:appendix-faq}

% ---------------------------------------------------------------------------
% Category 1: Foundational and Philosophical
% ---------------------------------------------------------------------------

\subsection{A complete causal model of a sociotechnical system might require thousands of variables. Does that mean causal inference is hopeless without a complete theory?}

No.
Science has never required complete theories, and causal inference is no exception.
The back-door criterion (Section~\ref{sec:primer-dags}) requires blocking all \emph{back-door paths} between treatment and outcome---not enumerating every variable in the universe.
A DAG with eight well-chosen nodes can achieve identification that a regression with two hundred covariates cannot, because what matters is whether the adjustment set blocks every confounding path, not how many variables appear in the model.
Moreover, sensitivity analysis methods allow researchers to quantify how strong an unmeasured confounder would need to be to overturn a finding~\citep{vanderweele2017sensitivity, oster2019unobservable} so that a study can provide credible evidence even without a complete causal model.
Meanwhile, design-based methods (Section~\ref{sec:primer-design}) sidestep the completeness problem entirely.
Difference-in-differences, panel fixed effects, and instrumental variables derive identification from features of the research setting---temporal variation, staggered adoption, quasi-random assignment---rather than from a claim that all confounders have been measured.
These methods require their own assumptions, but those assumptions concern the \emph{specific variation exploited}, not the completeness of the researcher's world model.

The broader philosophical stance is this: Every empirical study embeds causal assumptions, whether the researcher states them or not.
A regression that ``controls for'' five variables implicitly assumes that those five variables suffice to block all confounding---an assumption that is invisible, unexamined, and typically unjustified.
Causal inference does not demand that assumptions be \emph{correct}; it demands that they be \emph{transparent}.
Partial, explicit causal reasoning is strictly better than implicit, unexamined causal assumptions hiding inside every regression coefficient.

\subsection{My regression has $R^2 = 0.99$. Have I controlled for all confounders?}

No.
$R^2$ measures \emph{predictive} fit---how much variance in the outcome the model explains---not causal identification.
These two properties are logically independent, and conflating them is a common source of false confidence.

For example, suppose we want to estimate the effect of adopting code review ($T$) on post-release defect count ($Y$), and the true causal effect is $\beta_T = -5$ defects.
Developer inexperience ($U$) is an unmeasured confounder: Teams mandate code review for junior developers, and those same developers independently produce more defects.
We include lines of code (LOC) as a control variable, which is one of the strongest known predictors of defect counts in the empirical SE literature~\cite{DBLP:journals/tse/GravesKMS00, DBLP:journals/tse/EmamBGR01, DBLP:books/oreilly/11/HerraizH11}.
Regressing $Y$ on $T$ and LOC yields $R^2 \approx 0.99$ because LOC absorbs nearly all variance in defect counts.
But this near-perfect fit is causally misleading.
If LOC is a \emph{mediator}---code review encourages cleaner and smaller codebase that would have fewer defects---then conditioning on it blocks the indirect causal path $T \to \text{LOC} \to Y$, absorbing part of the genuine treatment effect into LOC's coefficient.
Meanwhile, the back-door path through $U$ remains wide open.
Because inexperienced developers both receive mandatory review ($U \to T$) and produce more defects ($U \to Y$), the unmeasured confounder induces a \emph{positive} spurious association between $T$ and $Y$---an upward bias.
Conditioning on LOC absorbs the genuine indirect effect while leaving this confounding bias intact.
If the confounding is strong enough, it can even \emph{flip the sign}.
A regression on $T$ and LOC may yield $\hat{\beta}_T \approx +1$, suggesting that code review \emph{increases} defects, when the true causal effect is $\beta_T = -5$.
The $R^2$ remains $0.99$ throughout, offering no warning that the coefficient is not merely biased but directionally wrong.
Had we instead omitted LOC and used a design-based strategy to address $U$, the model's $R^2$ might be a modest $0.35$---yet the causal estimate would be unbiased.

More generally, a model can achieve near-perfect $R^2$ while producing severely biased causal estimates.
This happens when the model includes \emph{mediators}---variables on the causal path from treatment to outcome---which block the genuine effect the researcher seeks to detect (Section~\ref{sec:primer-naive}).
It also happens when the model includes \emph{colliders}---common effects of treatment and outcome---which open spurious paths rather than closing them.
In both cases, $R^2$ rises because the added variable explains outcome variance, but the treatment coefficient moves \emph{away} from the true causal effect.
Conversely, a correctly specified causal model may have low $R^2$ because many factors affect the outcome but are unrelated to the treatment---the causal estimate can be unbiased even when the model explains little overall variance.

\subsection{I have a multivariate regression with 1000 covariates measuring everything I can think of in the study setting. Does that mean that I have sufficiently controlled for the important confounders?}

No---and in many cases, adding variables \emph{increases} bias.
This counterintuitive result follows directly from the causal structure of the data-generating process, which may make adding a variable introduce \emph{mediator bias} and \emph{collider bias} while not addressing unobserved confounders (see Section~\ref{sec:primer-naive} and Section~\ref{sec:primer-dags}).
Consequently, ``throwing in everything available'' is not a conservative strategy---it is a gamble whose expected payoff depends on the causal structure, which regression alone cannot reveal~\cite{wysocki2022statistical}.
The back-door criterion provides the principled alternative: Include variables that block confounding paths; exclude mediators, colliders, and their descendants.

\subsection{What is the role of DAGs in selection-on-observables vs.\ design-based identification?}

The two identification families use DAGs differently, and understanding this distinction clarifies what each requires from the researcher.
In \emph{selection-on-observables} methods---regression, matching, inverse probability weighting---DAGs are \emph{essential}.
The back-door criterion (Section~\ref{sec:primer-dags}) is the only principled basis for deciding which covariates to include and which to exclude.
Without a DAG, the researcher has no formal basis for determining whether the covariate set blocks all confounding paths, and covariate selection defaults to convenience, tradition, or data availability---none of which bear any necessary relationship to causal identification.
In \emph{design-based identification}---difference-in-differences, instrumental variables, regression discontinuity, panel fixed effects---DAGs play a \emph{supporting} role.
The design itself provides identification through quasi-random variation, so the researcher does not need a complete DAG to determine the adjustment set.
However, a partial DAG remains valuable for three purposes: clarifying which confounders the design absorbs (e.g., panel FE absorbs all time-invariant confounders), assessing whether design-specific assumptions are plausible (e.g., whether the exclusion restriction for IV holds, or whether unmeasured time-varying confounders threaten panel FE), and reasoning about the mechanisms through which the treatment affects the outcome.
One who draws even a partial DAG reasons more clearly about what their design does and does not address.
Both identification families benefit from making assumptions visual and debatable~\cite{imbens2020potential}.

% ---------------------------------------------------------------------------
% Category 2: "My Current Methods Work Fine"
% ---------------------------------------------------------------------------

\subsection{I already control for confounders in my regressions and report threats to validity. What does causal inference add?}

Two things: discipline and tools:
\begin{enumerate}[leftmargin=15pt]
    \item \emph{Discipline}:
    Causal inference forces the researcher to define the estimand before choosing a method (Section~\ref{sec:guide-question}), construct a DAG that makes every assumption explicit, and justify each covariate on causal---not statistical---grounds.
    This discipline routinely reveals problems invisible to the standard regression workflow: The study answers a different question than intended (e.g., ATT when the research question concerns ATE), a key covariate is a mediator whose inclusion blocks the effect of interest, or a ``control variable'' is a collider whose inclusion opens a spurious path.
    The estimand-first principle~\cite{lundberg2021estimand} ensures that the statistical procedure is evaluated against a well-defined target quantity, not assessed only on internal statistical criteria.
    \item \emph{Tools}:
    Design-based methods---difference-in-differences, panel fixed effects, instrumental variables---can absorb entire categories of unmeasured confounders that no regression can adjust for (Section~\ref{sec:primer-design}).
    Panel fixed effects, for example, remove \emph{all} time-invariant confounders---developer experience, team culture, organizational practices---without requiring them to be measured.
    These methods are strictly more powerful than covariate adjustment in settings where SE data has the temporal or panel structure that design-based identification exploits.
\end{enumerate}

\subsection{What is wrong with interpreting regression coefficients as causal effects?}

A regression coefficient has a causal interpretation \emph{only} when specific identifying assumptions are satisfied: conditional ignorability (all confounders measured), correct functional form, and no conditioning on mediators or colliders (Section~\ref{sec:primer-po}).
Without explicitly stating and defending these assumptions, the coefficient is a descriptive association whose causal content is unknown.
If a researcher simply fits a regression, reports coefficients, and interprets it in causal language in the discussion, this is exactly the pattern that the ``Table~2 fallacy''~\cite{westreich2013table2} describes.
Being careful about modeling choices (variable selection, outlier treatment, functional form) is necessary but not sufficient: It addresses statistical concerns while leaving the causal identification question unanswered.
A well-specified regression with all the right controls is a valid identification strategy---but the researcher must \emph{argue} that the controls are right, using a DAG or substantive domain knowledge, not hope that statistical diligence substitutes for causal reasoning.

\subsection{I run controlled experiments with developers. Why do I need observational causal inference methods?}

Controlled experiments are the gold standard for internal validity and nothing in this tutorial diminishes their value.
But many of the most important SE questions---Does adopting a language reduce defects in production?
Does code review improve long-term maintainability?
Does CI/CD adoption increase deployment frequency?---cannot be answered with lab experiments because they involve organizational decisions, long time horizons, and populations that cannot be randomly assigned.
Observational causal inference methods are designed precisely for these settings.

The target trial framework (Section~\ref{sec:primer-causality}) bridges the two worlds: It asks the researcher to describe the experiment they wish they could run and then evaluates how far their observational design departs from that ideal.
SE researchers who are already comfortable with controlled experiments~\cite{DBLP:books/sp/WohlinRHORW24} will find this framing natural---it simply extends the experimental logic to settings where randomization is infeasible.

\subsection{If I run a randomized experiment, am I free from causal reasoning?}

No.
Randomization is powerful but not a blanket exemption from causal reasoning.
Several common situations break or weaken the experimental guarantee, and all are prevalent in SE research, including:
\begin{enumerate}[leftmargin=15pt]
    \item \emph{Small samples.}
    Randomization balances confounders \emph{in expectation}, but with small $n$---common in developer studies with 20--40 participants---realized imbalance can be severe.
    The ``treatment'' and ``control'' groups may differ substantially on key characteristics by chance alone.
    As an SE-specific example, \citet{DBLP:conf/icse-seip/ParadisGMNMMZFC25} ran an RCT with 96 Google developers to measure AI's effect on coding speed; the unadjusted comparison was statistically significant ($p = .038$), but the estimate lost significance ($p = .086$) after controlling for developer seniority and daily coding hours (i.e., $X$), illustrating how even moderate sample sizes leave room for chance imbalance to shift conclusions.
    Checking and reporting covariate balance, and using stratified or blocked randomization, is essential~\cite{DBLP:books/sp/WohlinRHORW24}.
    \item \emph{Non-compliance.}
    If participants do not follow their assigned treatment---e.g., developers assigned to use a new tool revert to the old one---the naive comparison of assigned groups estimates the intention-to-treat effect, not the treatment effect itself.
    Recovering the latter requires instrumental variable reasoning, using the random assignment as an instrument to estimate the local average treatment effect for compliers~\cite{angrist1996identification}.
    \item \emph{Non-randomizable attributes.}
    Many SE-relevant ``treatments'' cannot be randomly assigned, e.g., developer experience, gender, educational background, programming language fluency.
    These are inherent characteristics, not interventions.
    Studying their effects requires the same observational identification strategies this tutorial covers; calling a group comparison an ``experiment'' does not change the identification problem.
    \item \emph{SUTVA violations.}
    If treated and control group subjects interact (e.g., developers may work on the same team, same codebase, same communication channels), the treatment can spill over, violating the stable unit treatment value assumption (Section~\ref{sec:primer-po}) that a randomized controlled experimental design requires.
    \item \emph{External validity.}
    A lab experiment with developers on short-term tasks may have clean identification but tell us nothing about production software at scale and in the long run.
    The target trial framework (Section~\ref{sec:primer-causality}) helps evaluate this gap between the study's actual scope and the research question's intended scope.
\end{enumerate}
In general, causal reasoning is not a tax imposed on observational studies that experiments get to avoid.
It is the discipline of thinking clearly about what your evidence can and cannot show, regardless of the study design.

\subsection{My paper was accepted at a top SE venue without any of this. Why should I change?}

Publication norms evolve.
In other empirical disciplines, growing awareness of the gap between causal ambition and methodological practice has led to concrete reforms---pre-registration, registered reports, estimand-first workflows, and editorial requirements for explicit identification strategies~\cite{grosz2020taboo, lederer2019control}.
SE has not yet adopted these norms, but the trajectory toward greater methodological rigor is clear.

Our analysis (Section~\ref{sec:bg-adoption}) shows that 29\% of empirical SE papers ask causal questions but fewer than 4\% use causal methods---a gap that reviewers and editors are increasingly aware of.
Similar expectations for explicit identification strategies will likely emerge in SE venues.

Early adoption is both a contribution to the field and a practical advantage: A study that explicitly states its identification strategy, constructs a DAG, and conducts sensitivity analysis is harder to criticize, more likely to produce durable findings, and more useful to practitioners and policymakers than one that relies on regression controls and a formulaic threats-to-validity section.

% ---------------------------------------------------------------------------
% Category 3: "SE Data Is Too Messy"
% ---------------------------------------------------------------------------

\subsection{SE data is noisy and poorly measured. How can any quasi-experimental design be credible here?}

SE data has genuine measurement challenges---noisy defect proxies, construct validity issues, potential SUTVA violations in interconnected ecosystems---and this tutorial acknowledges them honestly (Section~\ref{sec:guide-credibility}).
But SE data also has \emph{structural advantages} that many social science datasets lack.

Repositories provide fine-grained temporal data at the commit level.
Developers contribute across projects, creating within-person panel variation that panel fixed effects can exploit.
Tools and practices are adopted in staggered waves, creating natural difference-in-differences settings.
Platform policies---e.g., GitHub Actions rollout, npm security mandates, language deprecation announcements---create exogenous shocks that resemble natural experiments.

The question is not whether SE data is as clean as a minimum-wage natural experiment~\cite{card1994minimum} but whether the available variation, combined with domain knowledge and sensitivity analysis, supports a credible argument.
The pragmatic stance in Section~\ref{sec:guide} walks the researcher through this assessment, and the empirical standard in Appendix~\ref{sec:appendix-standard} provides verifiable criteria for reviewers.

\subsection{Parallel trends, exclusion restrictions, strict exogeneity---these assumptions are hard enough to defend in economics. Can SE researchers defend them credibly?}

This is a legitimate concern, and the tutorial does not claim that design-based identification is easy.
The pragmatic stance (Section~\ref{sec:guide}) exists precisely to walk SE researchers through the reasoning process, and the empirical standard (Appendix~\ref{sec:appendix-standard}) provides verifiable criteria for reviewers.

The pedagogical bet is that SE researchers---who are trained to reason about complex systems, dependencies, and graph structures---can learn to construct and defend identification arguments if given the right conceptual framework.
The alternative---continuing to make implicit causal claims from regression without any identification strategy---is not more credible; it is simply less transparent.
A DiD study that states the parallel trends assumption, tests it with pre-trend plots, and discusses remaining threats is epistemically superior to a regression study that makes causal claims without acknowledging any identification assumption at all.

\subsection{Aren't design-based assumptions just as untestable as the no-unmeasured-confounders assumption? What have you gained?}

Three things.

First, \emph{transparency}: Design-based assumptions are named, specific, and debatable---parallel trends, the exclusion restriction, continuity at the threshold.
The no-unmeasured-confounders assumption is a blanket claim that every relevant variable has been measured, a claim that is both stronger and harder to evaluate.

Second, \emph{partial testability}: Pre-trend tests for DiD~\cite{roth2023whats}, first-stage F-tests for IV, McCrary density tests for RDD, and overidentification tests provide empirical evidence---not proof---bearing on the assumptions.
No analogous diagnostic exists for the claim that all confounders are measured.

Third, \emph{domain-specificity}: The assumption shifts from ``I have measured everything relevant'' (which SE researchers almost never can defend from repository data) to ``the specific variation I exploit is quasi-random'' (which is a more tractable claim in settings with staggered adoption, platform shocks, or within-person variation).
The assumptions are not weaker---they are \emph{different}, and they are different in a way that aligns with the strengths of SE data.

\subsection{A clever identification strategy tells you about a local effect. Isn't that useless for practitioners?}

External validity is a real limitation of design-based methods, and the tutorial discusses the internal--external validity trade-off explicitly (Section~\ref{sec:primer-stance}).
A DiD study of TypeScript migration estimates the effect for projects that migrated, not for all projects.
An IV study recovers the LATE for compliers, not the ATE for the population.
These are genuinely narrower than the questions practitioners might ideally like answered.

But the alternative is not ``general causal knowledge from regression''---it is no causal knowledge at all, disguised as general knowledge.
A credible local estimate (e.g., the effect of TypeScript adoption on defects among projects that migrated) is more useful to a practitioner than a biased general estimate (e.g., the raw correlation between TypeScript use and defects across all GitHub projects), because the former has a causal interpretation while the latter conflates the treatment effect with selection bias.

Accumulating multiple credible local estimates across settings and populations---triangulation~\cite{shadish2002experimental}---is how fields build general causal knowledge.
This is precisely how economics built its evidence base on minimum wage effects, returns to education, and health insurance coverage: not through a single definitive study but through dozens of credible local estimates that converge.

% ---------------------------------------------------------------------------
% Category 4: Practical Guidance
% ---------------------------------------------------------------------------

\subsection{I am a new PhD student and my advisor wants me to study causal questions in SE. Where do I start?}

Start with the causal question, not the method.
Read Section~\ref{sec:primer} for the conceptual foundations, including the pragmatic stance in Section~\ref{sec:guide}.

Concretely: (1)~Articulate the target trial---what experiment would you run if you could?
(2)~Sketch a DAG encoding your beliefs about what causes what.
(3)~Look at your data through the lens of ``what quasi-experimental variation exists?''---panel structure, staggered adoption, threshold-based assignment.
(4)~Choose the method that matches the variation (Table~\ref{tab:design-features}).

For textbooks, start with \citet{cunningham2021causal} for an accessible introduction with code examples in R, Stata, and Python, then \citet{angrist2009mostly} for deeper treatment of the potential outcomes framework.
\citet{hernan2020causal} is essential for the target trial framework and the connection between causal questions and study design.
\citet{huntingtonklein2021effect} provides an excellent bridge between causal concepts and practical implementation.

\subsection{What software tools and packages should I use?}

The choice of tool matters far less than the choice of identification strategy---use whatever language your project already uses.

For \textbf{R}: \Code{fixest} for fast panel fixed effects, DiD, and IV estimation; \Code{did} for the \citet{callaway2021difference} staggered DiD estimator; \Code{rdrobust} for regression discontinuity; \Code{MatchIt} for matching; \Code{sensemakr} for sensitivity analysis and robustness values~\cite{cinelli2020making}; and \Code{dagitty} for DAG construction and identification analysis.
For \textbf{Python}: \Code{linearmodels} for panel FE and IV; \Code{dowhy} for end-to-end causal inference workflows; and \Code{causaldata} for textbook replication datasets.
For \textbf{Stata}: \Code{reghdfe} for high-dimensional fixed effects; \Code{did\_multiplegt} for staggered DiD; and \Code{rdrobust} for RDD\@.
For \textbf{DAG construction}: DAGitty (\url{https://dagitty.net}) provides a web-based interface for drawing DAGs and computing adjustment sets; TikZ in \LaTeX{} produces publication-quality figures (as in this paper).

\subsection{How do I construct a DAG? I don't know all the causal relationships in my domain.}

You do not need to know all causal relationships---you need to encode your \emph{assumptions} about the ones relevant to your research question.
The DAG is a model of the researcher's beliefs, not a claim of omniscience~\cite{pearl2009causality}.

Start with three nodes: treatment, outcome, and the most obvious confounder.
Then ask three questions iteratively.
``What else affects both treatment and outcome?''---add those as confounders.
``Does the treatment affect the outcome through intermediate variables?''---add those as mediators (and remember not to condition on them).
``Are there common effects of treatment and outcome?''---those are colliders (and must not be conditioned on).

Read prior literature to identify established relationships, talk to domain experts (e.g., developers, project managers) to surface tacit knowledge, and iterate.
A DAG is a living document, not a one-time exercise.
The value lies in the process of making assumptions explicit and debatable---not in getting the DAG ``right'' on the first try.
Figure~\ref{fig:pl-dag} provides a worked example of this process for the programming language and defect proneness question.

\subsection{My data is cross-sectional. Can I still do causal inference?}

Cross-sectional data---a single snapshot in time---severely limits causal identification.
Panel fixed effects (no temporal variation), difference-in-differences (no before/after comparison), and most design-based methods are unavailable.

The remaining options are:
(1)~Selection-on-observables (regression, matching, IPW) with a carefully justified DAG and sensitivity analysis~\cite{cinelli2020making}---but you must honestly acknowledge that the no-unmeasured-confounders assumption is strong and likely violated in most SE settings where key confounders are unmeasurable from repository data.
(2)~Instrumental variables, if you can find a plausible instrument.
(3)~Regression discontinuity, if the setting features a threshold-based assignment mechanism.

Often the most productive response is to \emph{reformulate the question} to exploit available variation (Section~\ref{sec:guide-neither}).
Can you collect panel data instead of cross-sectional data?
Can you identify migration or adoption events that create before/after variation?
The guide's most creative contribution is often in transforming an intractable cross-sectional question into a tractable longitudinal one---for example, ``Does language affect code quality?'' becomes ``Does adopting a type system reduce defects?'' when migration events are observable.

\subsection{How do I convince reviewers that my causal claims are credible?}

Follow the empirical standard in Appendix~\ref{sec:appendix-standard}.
At minimum: state the causal question explicitly, articulate the target trial, define the estimand, name the identification strategy, state the identifying assumptions, justify covariate selection using a DAG or substantive argument, conduct sensitivity analysis or robustness checks, and discuss at least one alternative explanation with evidence.

Use causal language (``X affects Y,'' ``X reduces Y'') only when your identification strategy supports it.
Use hedged language (``consistent with a causal interpretation under the assumption of\ldots'') when it does not---this is a strength, not a weakness (see the Invalid Criticisms section of the empirical standard).

A well-structured limitations section that \emph{mitigates} threats with quantitative tools---sensitivity analysis (how strong must an unmeasured confounder be to nullify the result?), falsification tests (does the treatment ``affect'' outcomes it should not?), robustness checks (does the result survive alternative specifications?)---is far more convincing than a laundry list of potential problems accompanied by no engagement with their severity.

% ---------------------------------------------------------------------------
% Category 5: Scope and Boundaries
% ---------------------------------------------------------------------------

\subsection{What about causal discovery algorithms that learn causal structure from data?}

Causal discovery algorithms---such as PC, GES, and FCI~\cite{spirtes2001causation}---attempt to learn DAG structure from observational data under assumptions like faithfulness and causal sufficiency.
They are a valid and active research direction, but they face practical challenges in SE settings.

\citet{journals/ese/HulseEM25} showed that causal graph generators applied to SE data are unstable, with over half of discovered edges changing across minor data perturbations.
The required assumptions---particularly faithfulness (every conditional independence in the data reflects a structural independence in the graph) and causal sufficiency (no unmeasured common causes)---are often implausible with noisy, high-dimensional SE data.

This tutorial focuses on theory-driven DAGs because they leverage the researcher's \emph{domain knowledge}---the comparative advantage of SE researchers who understand the systems they study---and because they are compatible with design-based identification, which does not require learning the full causal structure.
The two approaches are complementary: Data-driven methods can suggest edges to investigate, while domain knowledge adjudicates which edges are plausible and which are artifacts of data noise.

\subsection{What about machine learning methods for causal inference, like causal forests or double/debiased ML?}

These methods are important recent advances that the primer mentions (Section~\ref{sec:primer-toolkit}) but does not cover in depth.
Causal forests~\cite{athey2016recursive} estimate \emph{heterogeneous} treatment effects---how the effect varies across subgroups---and double/debiased machine learning~\cite{chernozhukov2018double} allows flexible, high-dimensional functional forms while maintaining valid statistical inference.

However, these methods do \emph{not} solve the identification problem.
They still require an identification strategy---selection-on-observables or design-based---and the same identifying assumptions.
A causal forest trained on observational data with unmeasured confounders produces biased heterogeneous effect estimates, not unbiased ones; double/debiased ML with endogenous treatment assignment produces consistent estimates of a non-causal quantity.
The methods are estimation tools, not identification tools.

For the modal SE researcher who has not yet mastered the basics of identification, learning difference-in-differences and panel fixed effects will yield more immediate returns than learning causal forests.
Once identification is understood, these ML-based methods become powerful extensions for exploring effect heterogeneity and relaxing functional form assumptions.

\subsection{How does this tutorial relate to A/B testing as practiced in industry?}

A/B testing is a randomized controlled experiment---the gold standard for internal validity.
Nothing in this tutorial replaces A/B testing when it is feasible.
The tutorial addresses the vast majority of SE questions where A/B testing is infeasible: studying the effects of language choice, process adoption, or organizational practices in real-world settings where random assignment is impractical or unethical.

The target trial framework (Section~\ref{sec:primer-causality}) bridges the two worlds: It asks ``what A/B test would I run if I could?'' and then evaluates how far the observational design departs from that ideal.
Industry practitioners who run A/B tests can also benefit from the causal reasoning framework when analyzing the observational data that surrounds their experiments---e.g., long-term effects after the experiment ends, spillover effects on untreated teams, or generalization from the experimental population to the broader user base.


\section{An Empirical Standard for Causal Inquiries in Software Engineering}
\label{sec:appendix-standard}

This appendix presents an empirical standard for SE studies making causal claims from observational data, written in the format of the ACM SIGSOFT Empirical Standards~\cite{DBLP:journals/corr/abs-2010-03525} so it can be directly incorporated into the ACM SIGSOFT Empirical Standards collection as a new standard.
Where the pragmatic stance in Section~\ref{sec:guide} conveys the reasoning, this standard provides the minimum requirements that authors, reviewers, and editors can verify.

\subsection{Application}

This standard applies to empirical SE studies that:
\begin{itemize}[leftmargin=15pt]
    \item investigate a causal research question (i.e., whether an intervention, practice, tool, or design decision affects an outcome),
    \item use observational data (not a randomized controlled experiment) to draw causal conclusions, and
    \item employ an identification strategy to argue that the estimate has a causal interpretation.
\end{itemize}
This standard does not apply to purely descriptive or predictive studies, randomized controlled experiments (see the Experiments standard), or studies that explicitly disclaim any causal intent.
If a study uses observational data and makes causal claims---even implicitly---this standard applies in addition to any method-specific standard (e.g., Repository Mining, Data Science, Longitudinal).

\subsection{Specific Attributes}

\paragraph{Essential.}
\begin{itemize}[leftmargin=15pt]
    \item Explicitly states a causal research question (not merely correlational or predictive).
    \item Articulates the target trial: describes the hypothetical experiment the observational analysis attempts to emulate, specifying treatment, population, outcome, and timing.
    \item Defines the causal estimand (e.g., ATE, ATT, LATE) using potential outcomes notation or an equivalent plain-language definition.
    \item Names the identification strategy used (e.g., selection-on-observables, difference-in-differences, instrumental variables, regression discontinuity, panel fixed effects).
    \item States the formal identifying assumptions required by the strategy (e.g., conditional ignorability, parallel trends, exclusion restriction, strict exogeneity).
    \item EITHER justifies covariate selection using a causal model (e.g., back-door criterion applied to a DAG) OR provides a substantive argument for why each included covariate blocks a confounding path and each excluded covariate is not a confounder.
    \item Discusses which confounders the design absorbs and which remain as potential threats.
    \item For selection-on-observables: conducts sensitivity analysis quantifying vulnerability to unmeasured confounding (e.g., Rosenbaum bounds, E-value, robustness value, coefficient stability).
    \item For DiD: tests and reports pre-treatment trends.
    \item For IV: reports first-stage strength (e.g., F-statistic) and argues exclusion restriction plausibility.
    \item For RDD: argues continuity at the threshold; reports a density test (e.g., McCrary) or equivalent.
    \item Limitations are not merely listed but \emph{mitigated} with quantitative tools (e.g., sensitivity analysis, robustness checks, falsification tests).
    \item At least one plausible alternative explanation is considered and addressed with evidence or argument.
    \item Discusses measurement validity: whether proxies adequately capture the constructs of interest.
    \item Causal conclusions are appropriately qualified given the strength of the identification strategy.
    \item Reports effect sizes, not just statistical significance.
    \item Distinguishes clearly between causal and associational claims in its language.
\end{itemize}

\paragraph{Desirable.}
\begin{itemize}[leftmargin=15pt]
    \item Presents a causal DAG (or equivalent causal model) encoding the assumed data-generating process.
    \item Discusses whether the treatment is well-defined or acknowledges compound treatment ambiguity.
    \item Discusses SUTVA (no interference between units) or justifies why interference is negligible.
    \item Discusses the internal--external validity trade-off: how the identification strategy constrains generalizability.
    \item Conducts robustness checks under alternative specifications (e.g., different control sets, functional forms, sample restrictions).
    \item Conducts falsification or placebo tests (e.g., testing for effects on outcomes the treatment should not affect).
    \item Discusses the direction and likely magnitude of remaining bias.
    \item Provides data and code for reproduction.
    \item For design-based methods: discusses the extent to which the quasi-experimental variation approximates random assignment.
\end{itemize}

\paragraph{Extraordinary.}
\begin{itemize}[leftmargin=15pt]
    \item Triangulates across multiple identification strategies (e.g., both selection-on-observables and design-based) on the same question.
    \item Conducts formal sensitivity analysis using multiple methods (e.g., both Rosenbaum bounds and robustness values).
    \item Pre-registers the identification strategy and analysis plan.
\end{itemize}

\subsection{General Quality Criteria}

Internal validity (strength of causal identification), construct validity (quality of treatment and outcome measurement), transparency of assumptions, and appropriate qualification of conclusions.
External validity is important but should not override internal validity for causal claims.

\subsection{Examples of Acceptable Deviations}

\begin{itemize}[leftmargin=15pt]
    \item A DAG is not presented if the identification strategy is design-based (e.g., DiD, RDD) and the design-specific assumptions are argued substantively, though a DAG is still desirable.
    \item Formal sensitivity analysis is not conducted if the study honestly qualifies its conclusions as associational or ``consistent with a causal interpretation under stated assumptions'' rather than claiming definitive causation.
    \item A formally defined estimand (ATE, ATT, LATE) is not provided if the study clearly describes in plain language what causal quantity is being targeted and for what population.
    \item Pre-treatment trend tests are not reported for DiD if the treatment is a one-time event with only one pre-treatment period, provided the limitation is acknowledged.
\end{itemize}

\subsection{Antipatterns}

\begin{itemize}[leftmargin=15pt]
    \item Interpreting regression coefficients as causal effects without stating identifying assumptions (the ``Table~2 fallacy''~\cite{westreich2013table2}).
    \item Including covariates without causal justification; ``throwing in everything available'' increases the risk of conditioning on mediators, colliders, or their descendants.
    \item Claiming to ``control for confounders'' using regression without arguing that all important confounders are measured.
    \item Acknowledging that ``correlation does not imply causation'' in the limitations section but writing the discussion and implications as if the findings are causal.
    \item Using causal language (``X affects Y,'' ``X leads to Y,'' ``the impact of X on Y'') while disclaiming any causal intent---the mismatch between language and method misleads readers.
    \item Listing threats to validity without engaging with their severity, direction of bias, or mitigation.
    \item Conducting sensitivity analysis but only reporting it when results are robust, suppressing cases where conclusions are fragile.
\end{itemize}

\subsection{Invalid Criticisms}

\begin{itemize}[leftmargin=15pt]
    \item The study does not use a randomized controlled experiment.
    The purpose of causal inference methods for observational data is to draw credible causal conclusions when randomization is infeasible.
    \item The study does not measure all possible confounders.
    No observational study can measure all confounders; what matters is whether the identification strategy and sensitivity analysis provide a credible argument.
    \item The study ``only'' estimates a local effect (LATE, effect at the threshold).
    Local effects are valid causal estimates; criticizing generalizability is appropriate but rejecting on this basis alone is not.
    \item The study's causal conclusions are ``too cautious'' or ``too hedged.''
    Appropriate qualification of causal claims is a strength, not a weakness.
    \item The effect size is small.
    Small effects can be practically important at scale and are not grounds for rejection if precisely estimated.
    \item The study uses an ``old'' or ``simple'' method (e.g., DiD instead of a newer estimator).
    What matters is whether the method's assumptions are credible in the specific context, not whether a more complex method exists.
\end{itemize}

\subsection{Suggested Readings}

\begin{itemize}[leftmargin=15pt]
    \item \citet{angrist2009mostly}: Textbook on design-based causal inference (potential outcomes, IV, DiD, RDD).
    \item \citet{cunningham2021causal}: Accessible introduction to causal inference with applications.
    \item \citet{hernan2020causal}: Target trial framework and causal inference for observational studies.
    \item \citet{pearl2009causality}: Graphical causal models, DAGs, and do-calculus.
    \item \citet{rosenbaum2002observational}: Sensitivity analysis for observational studies.
    \item \citet{lundberg2021estimand}: The estimand-first principle for empirical research.
    \item \citet{cinelli2020making}: Robustness values and sensitivity to unobserved confounding.
    \item \citet{roth2023whats}: Modern DiD methods and pre-trend testing.
    \item \citet{imbens2020potential}: Comparison of potential outcomes and DAG approaches.
    \item \citet{DBLP:journals/tosem/GrafVlachyW24}: IV tutorial for SE researchers.
\end{itemize}

\subsection{Exemplars}

\begin{itemize}[leftmargin=15pt]
    \item \citet{DBLP:conf/sigsoft/ChengMCJGKZ022}: Panel fixed effects analysis of code quality and developer productivity at Google.
    \item \citet{DBLP:conf/icse/FangLHV22}: Difference-in-differences analysis of social media promotion on open-source project outcomes.
    \item \citet{DBLP:conf/icse/ChenSSGT26}: Difference-in-differences analysis of core contributor disengagement on open-source workflows.
    \item \citet{DBLP:conf/msr/HeMAKV26}: Difference-in-differences analysis of LLM coding agents on development velocity and code quality.
    \item \citet{DBLP:journals/tosem/FuriaTF24}: Structural causal models applied to programming language and defect data.
\end{itemize}

